\chapter{Corrosión y degradación ambiental}\label{ch:06}
\objetivos{Relacionar ambiente y microestructura con el daño; diferenciar pérdida general de sección, daño localizado y crecimiento asistido por el medio.}
\section{La corrosión como sistema}
La corrosión electroquímica necesita reacciones anódicas y catódicas y un camino de transporte de carga. En una zona anódica se oxida el metal; en la catódica se consumen electrones mediante una reacción compatible con el medio. La distribución puede variar a escala microscópica y no requiere dos piezas distintas. La temperatura, la humedad, el oxígeno, los contaminantes, los depósitos y las diferencias de aireación cambian las velocidades locales.

Una película pasiva puede reducir la velocidad, pero su existencia no garantiza inmunidad. Una discontinuidad, un depósito o un medio que desestabilice la película puede concentrar ataque. En inspección se distingue la cantidad de material perdido de la profundidad y distribución de picaduras: una pérdida media pequeña puede ocultar un defecto local significativo.
\section{Pares galvánicos y geometría}
Al conectar materiales de distinto comportamiento electroquímico en un electrolito, puede acelerarse la disolución del menos noble bajo esas condiciones. La serie galvánica depende del medio y del estado superficial. No se utiliza una tabla genérica como orden universal de compatibilidad.

La relación de áreas importa: un ánodo pequeño conectado a una gran área catódica puede sufrir una densidad de corriente anódica elevada. El aislamiento eléctrico, el control del agua, los sellantes y los recubrimientos forman parte de la prevención. Un recubrimiento defectuoso sobre el ánodo puede, en determinadas configuraciones, concentrar el ataque en una discontinuidad; la protección se estudia como conjunto.
\ilustra{corrosion}{Celda galvánica conceptual: el camino metálico conduce electrones y el electrolito mantiene la continuidad iónica. Las flechas no predicen una velocidad de corrosión ni un único sentido para todas las especies iónicas.}
\section{Corrosión bajo tensión y fatiga-corrosión}
La corrosión bajo tensión requiere una combinación susceptible de material y estado metalúrgico, ambiente y tensión de tracción, que puede incluir tensiones residuales. No es simplemente sumar una pérdida de espesor a un esfuerzo. Pueden aparecer grietas transgranulares o intergranulares según sistema y mecanismo.

Un umbral de intensidad en un medio, a menudo denominado $K_{ISCC}$, depende de condiciones de ensayo y no se intercambia con $K_{IC}$. La fatiga-corrosión combina carga variable y ambiente; puede favorecer iniciación y crecimiento y alterar el significado práctico de una meseta S--N. La frecuencia y las permanencias cambian el tiempo disponible para interacción ambiental.

La fragilización por hidrógeno también depende de susceptibilidad, aporte de hidrógeno y tensiones. Puede relacionarse con procesos de fabricación o mantenimiento, además del ambiente de servicio. Una pieza con aspecto superficial poco alterado puede no estar libre de un mecanismo de fragilización.
\section{Oxidación y corrosión caliente}
A alta temperatura, la formación de óxidos puede proteger o degradar según adherencia, continuidad, estabilidad y transporte a través de la capa. En un régimen controlado por difusión se utiliza a veces una ley parabólica para el espesor de óxido $x$:
\begin{equation}x^2-x_0^2=k_p t.\label{eq:oxidacion}\end{equation}
La constante tiene unidades de longitud al cuadrado por tiempo y depende de temperatura y ambiente. También se usan leyes con ganancia de masa por área; sus constantes no se intercambian con las de espesor. Un régimen lineal, desprendimientos o una transición de mecanismo invalidan el uso indiscriminado de la ley parabólica.

La corrosión caliente puede involucrar depósitos salinos y degradación acelerada de la protección. No se identifica únicamente con «oxidación más rápida». Los recubrimientos de protección y las barreras térmicas tienen capas e interfaces, y pueden fallar por oxidación de la capa de enlace, crecimiento de óxidos, incompatibilidad termomecánica o erosión.
\caso{M06. Una constante necesita su definición}{En un ejercicio isotermo se supone $x_0=0$ y $k_p=4{,}0\times10^{-16}\,\mathrm{m^2/s}$. Tras $100\,\mathrm h$, $x=\sqrt{k_pt}=12{,}0\,\mu\mathrm m$. Tras $400\,\mathrm h$, el espesor es $24{,}0\,\mu\mathrm m$: cuadruplicar el tiempo duplica el espesor en esta ley. No se ha calculado pérdida de metal ni vida del recubrimiento, porque faltan densidades, estequiometría y criterio de fallo.
\par\textbf{Desarrollo para comparar tu intento.} La ley adoptada es $x^2=k_pt$ con espesor inicial cero. Aislar el espesor exige raíz positiva; antes de operar se convierten horas a segundos.

\begin{enumerate}

\item Convertir la primera duración.

\[\fitmath{t_{100}=100\cdot 3600\approx 360000\,\mathrm{s}}\]

\item Aplicar la raíz de producto constante por tiempo.

\[\fitmath{x_{100}=\sqrt{4\times10^{-16}\cdot 100\cdot 3600}\approx 1{,}2\times10^{-5}\,\mathrm{m}}\]

\item Repetir para cuatro veces el tiempo.

\[\fitmath{x_{400}=\sqrt{4\times10^{-16}\cdot 400\cdot 3600}\approx 2{,}4\times10^{-5}\,\mathrm{m}}\]

\end{enumerate}

}
\section{Inspección, prevención y registro}
La prevención combina elección de materiales, diseño de drenajes y uniones, protección superficial, limpieza compatible y control de contaminación. «Limpiar hasta que no se vea» no es un criterio de aceptación. Eliminar corrosión puede reducir sección y debe respetar datos aplicables de reparación y límites de daño.

Se registran ubicación, extensión, profundidad evaluada, método, incertidumbre, fotografías controladas y acciones. La repetición en una zona puede indicar un problema de acceso de humedad o un contacto galvánico no resuelto. En motores se vinculan hallazgos con ambiente de operación, almacenamiento y episodios de contaminación.

\section{La célula electroquímica y la geometría del ataque}
Una célula de corrosión necesita procesos anódico y catódico, conexión electrónica y un medio que permita transporte iónico. En una pareja galvánica importan potenciales en el ambiente real, polarización, área y resistencia del circuito. Un ánodo pequeño conectado a un cátodo grande puede sufrir una densidad de corriente anódica elevada. Invertir las áreas cambia la severidad aunque permanezca la misma pareja nominal.

Los huecos retienen electrolito y favorecen diferencias de aireación; un ataque localizado puede progresar sin una pérdida de masa total llamativa. Por ello la inspección de uniones, zonas de drenaje y superficies ocultas se relaciona con geometría y acceso. Un acabado protege mientras conserva continuidad y compatibilidad; una discontinuidad puede concentrar ataque en una zona pequeña.
\section{Daño ambiental que interactúa con tensiones}
La corrosión bajo tensión requiere una combinación susceptible de material y estado, ambiente y tensión de tracción, incluida la residual. No es simplemente corrosión general acelerada por cualquier carga. La fatiga--corrosión combina ciclos y ambiente y puede modificar iniciación, umbral y crecimiento; las curvas en aire no se trasladan sin más a ese entorno. En ambos casos es esencial identificar mecanismo y condiciones, no solo medir una dimensión final del defecto.

La fragilización por hidrógeno puede relacionarse con procesos de fabricación, protección o servicio, según material y exposición. Un tratamiento o protección que mejore una característica no garantiza ausencia de efectos secundarios. El análisis de fallo examina procedencia, microestructura, superficie de fractura y secuencia temporal; una sola fotografía raramente prueba la causa.
\section{Oxidación protectora y pérdida de protección}
Una película compacta y adherente puede frenar difusión y dar una cinética aproximadamente parabólica. Una película porosa, fisurada o que se desprende puede no mantener esa protección. Los ciclos térmicos introducen incompatibilidades de dilatación y tensiones; el espesor de óxido no equivale directamente a espesor de metal consumido sin conocer composición, densidades y geometría.

En zonas calientes también puede haber corrosión por depósitos y especies ambientales que alteran la película. No se agrupan todos los ataques de alta temperatura bajo un único parámetro de oxidación. En mantenimiento, limpieza e inspección deben ajustarse a datos del producto: una retirada agresiva de depósito puede dañar el sustrato o borrar evidencia.
\repaso{¿Una unión carbono--aluminio es segura frente a corrosión por no contener dos metales?}{No puede concluirse así. El carbono conductor puede participar en un acoplamiento electroquímico con el aluminio en un ambiente adecuado. Se estudian aislamiento, protección, sellado, drenaje y mantenimiento dentro del diseño aplicable.}

\resumen{La corrosión no se caracteriza solo por masa perdida. La geometría y la conexión entre materiales importan. Una ley de oxidación especifica qué magnitud ajusta y bajo qué régimen. $K_{ISCC}$, la resistencia estática y los datos de fatiga en aire no son criterios intercambiables.}
\fuentes{FAA--General, cap. 8, limpieza y control de corrosión \cite{FAA-G}; introducción docente a materiales; protección y materiales de motores en NASA \cite{NASA-CMC}.}
